\documentclass[instructions.tex]{subfiles}
\begin{document}
 
\chapter{De la confession sacrilége.}
 
\primo On abuse du sacrement toutes les fois qu’on y cache avec connoissance un péché mortel, ou lorsqu’on ne veut pas quitter un péché mortel qu’on accuse. Les péchés ne sont pas remis par une telle confession, et on se rend coupable d’un nouveau crime, qui est le sacrilége; crime ordinairement plus énorme que les péchés qu’on accuse. C’est y ajouter un sacrilége encore plus horrible, lorsqu’on communie en cet état.
 
Vous découvririez les plaies de votre corps à un médecin, pour vous conserver la vie : pourquoi avez-vous honte de découvrir les plaies de votre conscience pour sauver votre âme ? Peut-on être guéri quand on ne veut pas déclarer son mal ? Il faut indispensablement au moins confesser ses péchés mortels, quand on le peut, pour en avoir le pardon, ou brûler pendant l’éternité : choisissez.
 
\secundo Prenez donc garde de vous laisser surprendre par la honte. Lorsqu’on y succombe, on s’accoutume insensiblement aux sacriléges; on s’en fait une habitude; cette habitude sacrilége étant formée, que doit-on espérer de celui qui abuse ainsi des remèdes mêmes du salut ? Il n’est point d’état plus malheureux et plus cruel. Les remords de la conscience, qui suivent le sacrilége, sont comme un bourreau qui ne donne point de relâche. Le démon profite de cet état pour porter une âme au découragement; et, pour la tromper, il lui persuade qu’il n’y a plus de ressource. Dans ces troubles et ces agitations, on s’ennuie de la vie; on perd la confiance, on se livre à toutes sortes de vices (1). Mais d’un autre côté, lorsqu’on envisage la mort, les jugemens de Dieu, l’éternité, on se sent saisi d’horreur. Oh ! quel état que celui d’une conscience ainsi tourmentée ! Quel repos peut avoir une âme qui porte ainsi, par sa malice, son supplice et son enfer avec elle ?
 
Le remède à de si grands maux, c’est d’avouer ses fautes et ses sacriléges avec confiance et avec humilité. Ne vaut-il pas mieux, dit saint Augustin, déclarer ses péchés, et subir une confusion d’un moment, devant un homme rempli de charité, que de souffrir une confusion accablante au jugement de Dieu, devant les anges et devant tout l’univers ?
 
\section{Résolutions}
 
\primo J’aurai toujours horreur de la confession sacrilége.
 
\secundo Pour l’éviter, je porterai constamment au tribunal sacré la sincérité la plus entière, afin de ne rien taire de tout ce que je dois y accuser; et une douleur amère de mes péchés, avec la résolution ferme de ne plus les commettre.
 
\tertio Enfin, si j’ai fait quelques confessions qui n’aient pas été entières par ma faute, ou si j’y ai manqué de contrition, je communiquerai avec franchise toutes mes peines à cet égard à mon confesseur, prêt à réparer le passé, suivant ses sages avis.

\prayer{Mon Dieu, pardonnez-moi l’abus que j’ai pu faire jusqu’ici d’un sacrement où vous faites éclater d’une manière si touchante votre miséricorde envers nous.}
 
\end{document}
